\section{Lifting the remaining cross-cup axes at coefficient two}
\label{sec:coefficient-two-cross-cups}

Write \(G=\widehat{\mathrm{GT}}\), \(H=G_{\mathbf Q}\),
\(M=\ker\lambda\), and \(N=\ker(\lambda,\rho_2^I)\).
The arithmetic field fixed by \(N\cap H\) is
\[
 K_2=\mathbf Q(\mu_\infty,2^{1/\infty}),\qquad
 A=\operatorname{Gal}(K_2/\mathbf Q)=G/N.
\]
For an odd prime \(p\), let \(c_p\) be the quadratic
character of \(p^*=(-1)^{(p-1)/2}p\), and let
\(\delta_p\) be the degree-two generator of the cyclic
two-primary quotient of \(\mathbf Z_p^\times\), inflated
to the cyclotomic group and then to \(A\). Put
\[
 \begin{gathered}
 b(u)=\frac{u_2-1}{2}\pmod2,\qquad
 a(u)=\frac{u_2^2-1}{8}\pmod2,\\
 \delta_2=b^2,\qquad A_p=\frac{p^2-1}{8}\pmod2.
 \end{gathered}
\]
All cohomology in this section has trivial coefficient
module \(\mathbf F_2\).

\begin{theorem}\label{thm:coefficient-two-cross-cups}
For every odd prime \(p\), there is a continuous
\(G\)-invariant homomorphism
\[
                     \xi_{2,p}:M\longrightarrow\mathbf F_2
\]
whose restriction to \(N\) transgresses to
\begin{equation}
 e_{a,c_p}=a c_p+A_p(\delta_2+\delta_p)
                        \quad\hbox{in }H^2(A,\mathbf F_2).
                                                        \label{eq:coefficient-two-cross-1}
\end{equation}
In particular \eqref{eq:coefficient-two-cross-1} has zero inflation to \(G\), and
\begin{equation}
                       \xi_{2,p}|_{G_{K_2}}=\kappa_{2,p},
                                                        \label{eq:coefficient-two-cross-2}
\end{equation}
where \(\kappa_{2,p}\) is the unique arithmetic target
character transgressing to \eqref{eq:coefficient-two-cross-1}.

The character is obtained from the action on compact
homology, modulo eight, of the degree-\(p\) cyclic cover
of the three-punctured sphere with monodromies
\(1,1,-2\). The same construction is defined and invariant
on the multiplier-one subgroup of any group of continuous
positive peripheral automorphisms of \(\widehat F(x,y)\).
The transgression assertion above concerns specifically
the full quotient \(G/N=A\).
\end{theorem}

The proof also constructs a companion character whose
arithmetic restriction is \(\psi_p\). Its correction
terms provide a useful check on the normalization.

\subsection{Integral homology and the multiplier}

Set
\[
 U_p=\ker\bigl(\widehat F(x,y)\longrightarrow C_p,
                         \quad x,y\longmapsto1\bigr),
 \qquad m=p-1.
\]
The cover has three punctures and compact genus
\((p-1)/2\). Let \(V\) be its first homology over
\(\mathbf Z_2\), modulo the peripheral homology submodule.
It is free of rank \(m\), with its unimodular alternating
intersection form. The deck generator \(D\) has each
nontrivial \(p\)-th-root character once.
Indeed, on nontrivial character summands the cellular
boundary map
\[
 \mathbf Z_2[C_p]^2\longrightarrow\mathbf Z_2[C_p],
          \qquad (r,s)\longmapsto(D-1)(r+s)
\]
has a rank-one kernel. Its two trivial summands are
exactly the peripheral homology. The latter assertion
is integral: \(p\) is a unit in \(\mathbf Z_2\).

A positive peripheral automorphism with common multiplier
\(\lambda\) preserves \(U_p\) and gives an integral
similitude
\begin{equation}
 \rho_{p,2}(g)\in\operatorname{GSp}(V),\qquad
 \rho_{p,2}(g)D\rho_{p,2}(g)^{-1}
       =D^{\lambda_g\bmod p},\qquad
 \operatorname{sim}(\rho_{p,2}(g))=\lambda_{g,2}.
                                                        \label{eq:coefficient-two-cross-3}
\end{equation}
We verify the multiplier at two without passing to an
exponent-two group. A surface presentation is
\begin{equation}
 \prod_{j=1}^{(p-1)/2}[a_j,b_j]c_1c_2c_3=1.
                                                        \label{eq:coefficient-two-cross-4}
\end{equation}
There is one puncture above each original puncture.
The conjugator of a peripheral power can therefore be
moved into \(U_p\) by multiplication by a power of the
corresponding original peripheral generator. Consequently
the automorphism sends each \(c_i\) to a
\(U_p\)-conjugate of \(c_i^\lambda\).

Take the torsion-free class-two pro-two quotient of the
free group on \(a_j,b_j,c_1,c_2\). Its collected
commutator coordinates form a free \(\mathbf Z_2\)-module.
After tensoring with \(\mathbf Q_2\), use the class-two
logarithm and project degree-two coordinates to
\(\bigwedge^2 V\). The logarithms of \(c_1,c_2\), and
their images, have zero projected central coordinate:
conjugation adds a wedge containing a peripheral class.
Equation \eqref{eq:coefficient-two-cross-4} gives projected central coordinate
\(-\sum_j a_j\wedge b_j\) for \(\log c_3\).
The image of \(c_3\) multiplies that coordinate by
\(\lambda_2\). Thus
\[
     (\bigwedge^2\rho_{p,2}(g))
          \left(\sum_j a_j\wedge b_j\right)
       =\lambda_{g,2}\sum_j a_j\wedge b_j.
\]
Both sides are integral. This is the alternating
similitude identity, proving \eqref{eq:coefficient-two-cross-3}. This argument uses
only the stated peripheral action and the presentation
of the cover.

Let \(E\) be the normalizer of \(\langle D\rangle\)
in \(\operatorname{GSp}(V)\). Its quotient map records
deck exponent and similitude multiplier:
\begin{equation}
 \pi:E\longrightarrow Q_p,
       \qquad Q_p=C_m\times\mathbf Z_2^\times.
                                                        \label{eq:coefficient-two-cross-5}
\end{equation}
We construct two \(E\)-invariant characters of
\(T=\ker\pi\), denoted \(\chi_a,\chi_b\), and compute
their central extensions exactly.

\subsection{A norm-one coinvariant calculation}

\begin{lemma}\label{lem:two-norm-one-coinvariants}
Let \(L/\mathbf Q_2\) be unramified of degree \(2s\),
with integer ring \(O\), arithmetic Frobenius \(\sigma\),
and involution \(\bar z=\sigma^s(z)\). Let \(O_0\)
be the fixed integer ring, and put
\[
             T_s=\{z\in O^\times:z\bar z=1\}.
\]
There is a canonical isomorphism of finite groups
\begin{equation}
 \mathcal D_s:
 \mathbf Z_2^\times/(\mathbf Z_2^\times)^2
       \xrightarrow{\ \sim\ } T_s/(\sigma-1)T_s,
 \qquad
 t\longmapsto\left[\frac{\sigma(v)}v\right],
 \qquad v\bar v=t.
                                                        \label{eq:coefficient-two-cross-6}
\end{equation}
Here \((\sigma-1)T_s=\{\sigma(z)/z:z\in T_s\}\).
Define
\[
        \alpha_s=a\circ\mathcal D_s^{-1},\qquad
        \beta_s=b\circ\mathcal D_s^{-1}.
\]
For \(z\in T_s\cap(1+4O)\),
\begin{equation}
 \alpha_s(z)=
 \operatorname{Tr}_{\mathbf F_{2^s}/\mathbf F_2}
       \left(\frac{\log z}{4}\bmod2\right),
 \qquad \beta_s(z)=0.                                   \label{eq:coefficient-two-cross-7}
\end{equation}
Moreover
\begin{equation}
                  \alpha_s(-1)=1,
       \qquad \beta_s(-1)=s\pmod2.                       \label{eq:coefficient-two-cross-8}
\end{equation}
If \(v\bar v=-1\), then
\begin{equation}
                   \alpha_s(v^2)=1,
              \qquad \beta_s(v^2)=0.                    \label{eq:coefficient-two-cross-9}
\end{equation}
In \eqref{eq:coefficient-two-cross-9}, the square belongs to \(T_s\); the values are
independent of the choice of \(v\).
\end{lemma}

\begin{proof}
The norm on the units of an unramified extension is
surjective, and cyclic norm-one elements are Hilbert
90 quotients; see \cite[II, Corollary~1.23, p.~67;
III, Proposition~1.2, p.~98]{MilneCFT}.
For completeness, surjectivity on units follows by
first lifting the surjective residue-field norm and
then successively correcting each principal-unit
level using the surjective residue-field trace.

Choose \(v\) with \(v\bar v=t\). Replacing \(v\) by
\(vz\), \(z\in T_s\), changes its quotient in \eqref{eq:coefficient-two-cross-6}
by \(\sigma(z)/z\). Products of choices show that \eqref{eq:coefficient-two-cross-6}
is a homomorphism. A square \(t=w^2\), with
\(w\in\mathbf Z_2^\times\), has the choice \(v=w\),
so squares map to one.

Conversely, let \(z\in T_s\). Its full norm to
\(\mathbf Q_2\) is one. Hilbert 90 gives
\(z=\sigma(v)/v\) with \(v\in L^\times\). Multiplying
by a rational power of two makes \(v\) a unit because
the extension is unramified. The relative norm
\(v\bar v\) is fixed by \(\sigma\), hence lies in
\(\mathbf Z_2^\times\). This proves surjectivity.
If \(\sigma(v)/v=\sigma(z)/z\) for \(z\in T_s\),
then \(v/z\) is a rational unit, so \(v\bar v\) is
a square. This proves injectivity. The image of
\(\sigma-1\) is closed by compactness; the quotient
and \eqref{eq:coefficient-two-cross-6} are continuous.

To prove \eqref{eq:coefficient-two-cross-7}, choose \(v\) with
\(z=\sigma(v)/v\). Since \(z\equiv1\pmod4\), the
class of \(v\) modulo four is fixed by Frobenius.
Indeed its residue lies in \(\mathbf F_2^\times=\{1\}\);
write \(v=1+2y\pmod4\). Frobenius invariance then gives
\(y^2=y\pmod2\), so \(v\equiv1\) or \(3\pmod4\). Scaling by a rational
unit therefore makes \(v\equiv1\pmod4\).
Write \(v=\exp(4x)\). Then
\[
 t=v\bar v=\exp(4(x+\bar x))\in1+4\mathbf Z_2,
 \qquad \frac{\log z}{4}=\sigma x-x.
\]
Thus \(b(t)=0\) and \(a(t)=x+\bar x\pmod2\).
The reduction of \(\sigma x-x\) is \(x^2+x\), and
it belongs to \(\mathbf F_{2^s}\) because its relative
trace is zero. Its absolute trace telescopes:
\[
 \operatorname{Tr}_{\mathbf F_{2^s}/\mathbf F_2}(x^2+x)
                       =x^{2^s}+x=\bar x+x.
\]
This proves \eqref{eq:coefficient-two-cross-7}.

For \eqref{eq:coefficient-two-cross-8}, the unramified quadratic subfield contains
\(\delta=1+2\omega\), where \(\omega^2+\omega+1=0\).
It satisfies \(\delta^2=-3\) and
\(\sigma\delta=-\delta\). Its relative norm is
\(3\) for odd \(s\), and \(-3\) for even \(s\).
Hence \(\mathcal D_s(3)=[-1]\) or
\(\mathcal D_s(-3)=[-1]\), respectively.
Now \(a(3)=a(-3)=1\), while
\(b(3)=1\) and \(b(-3)=0\).

For \eqref{eq:coefficient-two-cross-9}, remove the odd-order Teichm\"uller factor of
\(v\). Its relative norm is one, so this operation
preserves the norm \(-1\); its square has zero value
under either quadratic character. We may thus take
\(v=1+2x\). The norm equation gives
\[
                    x+\bar x=1\pmod2,
       \qquad v^2=1+4(x+x^2)\in T_s\cap(1+4O).
\]
Formula \eqref{eq:coefficient-two-cross-7} and the same telescoping trace give
\(\alpha_s(v^2)=1\) and \(\beta_s(v^2)=0\).
Replacing \(v\) by another norm-minus-one unit
multiplies \(v^2\) by a square in \(T_s\), so the
values are independent of the choice.
\end{proof}

There is also a direct unit formula for the inverse square class.
If \(z=w/\bar w\in T_s\), then
\begin{equation}\label{eq:two-hilbert-norm-squareclass}
 \mathcal D_s^{-1}([z])=[N_{L/\mathbf Q_2}(w)].
\end{equation}
Changing \(w\) by a unit of the fixed field changes its absolute
norm by a rational square, so the right side is well-defined.
To prove the formula, write \(z=\sigma(v)/v\) and
\(v\bar v=t\in\mathbf Z_2^\times\), as above. Put
\[
 P=\prod_{j=0}^{s-1}\sigma^j(w),\qquad Z=vP.
\]
Then \(\sigma(Z)/Z=(\sigma(v)/v)(\bar w/w)=1\), so
\(Z\) is a rational unit. Therefore
\[
 N_{L/\mathbf Q_2}(w)=P\bar P=Z^2/t.
\]
The square classes of \(t^{-1}\) and \(t\) coincide, proving
\eqref{eq:two-hilbert-norm-squareclass}. In particular
\[
 \alpha_s(z)=a(N_{L/\mathbf Q_2}w),\qquad
 \beta_s(z)=b(N_{L/\mathbf Q_2}w).
\]
These formulas agree on the entire torus, including its first
principal-unit layer.

\subsection{Even Frobenius order}

Suppose \(d=\operatorname{ord}_p(2)=2s\), and let
\(r=m/d\). Every Frobenius orbit is self-dual. With
\(O\) as in Lemma~\ref{lem:two-norm-one-coinvariants},
we may write \(V=O^r\). On each summand the form is
\[
       (z,w)\longmapsto
       \operatorname{Tr}_{L/\mathbf Q_2}(d_i z\bar w),
       \qquad \bar d_i=-d_i,\quad d_i\in O^\times.
\]
The coefficients are units by unimodularity. Their
ratios lie in \(O_0^\times\); surjectivity of the
relative norm permits a rescaling making all of
them equal. Consequently
\begin{equation}
 T=T_s^r,
 \qquad \chi_a(z_1,\ldots,z_r)=\sum_i\alpha_s(z_i),
 \qquad \chi_b(z_1,\ldots,z_r)=\sum_i\beta_s(z_i).
                                                        \label{eq:coefficient-two-cross-10}
\end{equation}

Fix a primitive residue \(u\bmod p\). Choose the
Frobenius-orbit representatives
\(1,u,\ldots,u^{r-1}\). A lift of deck exponent
\(u\) cycles these summands, with Frobenius twist
\(\sigma^k\) on the last passage, where
\[
                  u^r=2^k\pmod p,
              \qquad \gcd(k,d)=1.
\]
In particular \(k\) is odd. The ratio between the
form coefficients before and after this last twist
lies in \(O_0^\times\); multiplying that summand by
a unit of prescribed relative norm restores the
form. This gives an integral symplectic lift \(J\).

For every \(t\in\mathbf Z_2^\times\), choose a unit
\(v\in O^\times\) with \(v\bar v=t\). Uniform scalar
multiplication by \(v\) gives a lift \(L(t)\) of
\((1,t)\) in \eqref{eq:coefficient-two-cross-5}. Thus \eqref{eq:coefficient-two-cross-5} is onto. These scalar
lifts commute with one another and with \(T\).
Conjugation by \(J\) only permutes factors and applies
Frobenius powers. Formula \eqref{eq:coefficient-two-cross-6} makes both characters
in \eqref{eq:coefficient-two-cross-10} invariant. Hence both are invariant under
all of \(E\).

The scalar correction used to define \(J\) commutes
with \(L(t)\). Their commutator has one nontrivial
component, \(\sigma^k(v)/v\), or its inverse according
to the choice of permutation convention. In the
coinvariant quotient of \eqref{eq:coefficient-two-cross-6}, this component is
\(\mathcal D_s(t)^k=\mathcal D_s(t)\). Therefore
\begin{equation}
             \chi_a([J,L(t)])=a(t),\qquad
             \chi_b([J,L(t)])=b(t).                       \label{eq:coefficient-two-cross-11}
\end{equation}

Over an algebraic closure, \(J\) cycles all \(m\)
deck eigenlines in one cycle. The product of its
scalar entries is \(-1\): the cycle has sign \(-1\)
and \(\det J=1\). Thus, exactly,
\begin{equation}
                             J^m=-I.                    \label{eq:coefficient-two-cross-12}
\end{equation}
Using \eqref{eq:coefficient-two-cross-8},
\begin{equation}
                    \chi_a(J^m)=r\pmod2,
           \qquad \chi_b(J^m)=rs\pmod2.
                                                        \label{eq:coefficient-two-cross-13}
\end{equation}
For the cyclic subgroup generated by the multiplier
\(-1\), \eqref{eq:coefficient-two-cross-9} gives
\begin{equation}
                   \chi_a(L(-1)^2)=r\pmod2,
              \qquad\chi_b(L(-1)^2)=0.                  \label{eq:coefficient-two-cross-14}
\end{equation}
In particular the square of the lift of \(-1\) cannot
be discarded in this calculation.

\subsection{Odd Frobenius order}

Suppose \(d=\operatorname{ord}_p(2)\) is odd, and let
\(r=m/(2d)\). Let \(O\) be the integer ring in the
unramified extension of \(\mathbf Q_2\) of degree
\(d\). Frobenius orbits occur in distinct dual pairs.
The lattice is a sum of \(r\) hyperbolic \(O\)-pairs,
with form
\[
 \sum_i\operatorname{Tr}_{O/\mathbf Z_2}
                         (x_i y_i'-y_i x_i').
\]
Its torus and two characters are
\begin{equation}
 \begin{split}
 T&=(O^\times)^r,
       \quad z_i:(x_i,y_i)\longmapsto(z_i x_i,z_i^{-1}y_i),\\
 \chi_a(z)&=\sum_i a(N_{O/\mathbf Z_2}z_i),
       \qquad
 \chi_b(z)=\sum_i b(N_{O/\mathbf Z_2}z_i).
 \end{split}                                             \label{eq:coefficient-two-cross-15}
\end{equation}
They are invariant under Frobenius, inversion, and
permutation of pairs.

For a primitive residue \(u\), the quotient by
\(\langle2,-1\rangle\) has order \(r\), and
\(u^r=-2^k\) for some \(k\). A symplectic lift \(J\)
cycles the \(r\) pairs and on the final passage
exchanges reciprocal coordinates with a Frobenius
twist and one minus sign. A multiplicative similitude
section is
\[
           L(t):(x_i,y_i)\longmapsto(x_i,t y_i)
              \qquad(t\in\mathbf Z_2^\times).
\]
This proves surjectivity of \eqref{eq:coefficient-two-cross-5} and invariance of \eqref{eq:coefficient-two-cross-15}.
The commutator of \(J\) with \(L(t)\) has torus
parameter \(t\) or \(t^{-1}\) on one pair and one
elsewhere. Its norm is \(t^{\pm d}\). Since \(d\)
is odd, \eqref{eq:coefficient-two-cross-11} holds again. The eigenline-cycle proof
of \eqref{eq:coefficient-two-cross-12} also applies, and now
\begin{equation}
               \chi_a(J^m)=0,\qquad \chi_b(J^m)=r\pmod2.
                                                        \label{eq:coefficient-two-cross-16}
\end{equation}
The multiplicative section gives
\begin{equation}
                     \chi_a(L(-1)^2)
                         =\chi_b(L(-1)^2)=0.              \label{eq:coefficient-two-cross-17}
\end{equation}

\subsection{All central-extension coefficients}

For \(\chi=\chi_a\) or \(\chi_b\), the quotient
\begin{equation}
              1\longrightarrow C_2\longrightarrow
                 E/\ker\chi\longrightarrow Q_p
                    \longrightarrow1                   \label{eq:coefficient-two-cross-18}
\end{equation}
is a continuous central extension. Both characters
are nonzero: this follows from \eqref{eq:coefficient-two-cross-6} in even degree
and surjectivity of the unramified norm in odd degree.

The degree-two cohomology of
\(Q_p=C_m\times\mathbf Z_2^\times\), with
\(\mathbf Z_2^\times=\langle-1\rangle\times
\overline{\langle5\rangle}\), has basis
\begin{equation}
                 \delta_p,\quad b^2,\quad ba,
                         \quad b c_p,\quad a c_p.
                                                        \label{eq:coefficient-two-cross-19}
\end{equation}
This follows from the cyclic cohomology calculation
and the product formula; the procyclic factor
\(\overline{\langle5\rangle}\simeq\mathbf Z_2\)
has cohomological dimension one.

Equation \eqref{eq:coefficient-two-cross-11} gives mixed term \(a c_p\) for
\(\chi_a\) and \(b c_p\) for \(\chi_b\).
The preimage in \(E\) of the multiplier factor
\(\mathbf Z_2^\times\) is abelian in both descriptions,
so the coefficient of \(ba\) is zero.
Equations \eqref{eq:coefficient-two-cross-13}--\eqref{eq:coefficient-two-cross-17} give every other coefficient:
\begin{equation}
 \begin{array}{c|cc|cc}
 &\multicolumn{2}{c|}{\chi_a}
 &\multicolumn{2}{c}{\chi_b}\\
 &\delta_p&b^2&\delta_p&b^2\\ \hline
 d=2s &r&r&rs&0\\
 d\text{ odd}&0&0&r&0
 \end{array}                                             \label{eq:coefficient-two-cross-20}
\end{equation}
All entries are reduced modulo two. In both cases
the \(\chi_b\) entry at \(\delta_p\) is
\((p-1)/2\pmod2\).

The \(\chi_a\) entries equal \(A_p\). Indeed, if
\(d\) is odd, two is a square modulo \(p\), so
\(A_p=0\). If \(d\) is even, write
\(2=u^{rh}\), with \(\gcd(h,d)=1\); then \(h\)
is odd and the Legendre symbol \((2/p)=(-1)^r\).
The quadratic reciprocity supplement identifies
this sign with \((-1)^{A_p}\).
Consequently the exact extension classes are
\begin{equation}
 \begin{split}
 [E/\ker\chi_a]&=a c_p+A_p(b^2+\delta_p),\\
 [E/\ker\chi_b]&=b c_p+\frac{p-1}{2}\delta_p.
 \end{split}                                             \label{eq:coefficient-two-cross-21}
\end{equation}

These are exactly the arithmetic Hilbert corrections.
For \((2,p^*)\), the only possible finite ramification
places are \(2,p\), and both Hilbert signs are
\((-1)^{A_p}\). For \((-1,p^*)\), the only finite
sign is at \(p\), with exponent \((p-1)/2\).
The dyadic assertion uses \(p^*\equiv1\pmod4\).
These also follow directly from the explicit dyadic
Hilbert formula in \cite[VIII, \S5, p.~248]{MilneCFT}.
No coefficient in \eqref{eq:coefficient-two-cross-21} is inferred merely from
arithmetic Brauer vanishing.

\subsection{Characters, transgression, and the full target}

For \(g\in M\), \eqref{eq:coefficient-two-cross-3} puts \(\rho_{p,2}(g)\) in
\(T\). Define
\begin{equation}
               \xi_{2,p}(g)=\chi_a(\rho_{p,2}(g)),
       \qquad \widetilde\psi_p(g)=\chi_b(\rho_{p,2}(g)).
                                                        \label{eq:coefficient-two-cross-22}
\end{equation}
They are homomorphisms, and invariance of the torus
characters makes them invariant under conjugation
by every positive peripheral automorphism in the
ambient group. They are continuous. More precisely,
\eqref{eq:coefficient-two-cross-7} kills \(T_s\cap(1+8O)\) for \(\alpha_s\) and
\(T_s\cap(1+4O)\) for \(\beta_s\). In odd degree
the same assertions follow from the norm formulas.
Thus the two characters in \eqref{eq:coefficient-two-cross-22} depend only on
compact homology modulo eight and four, respectively.

The representation \(G\to E\), followed by
either quotient in \eqref{eq:coefficient-two-cross-18}, lifts the cyclotomic map
\(G\to Q_p\). Both classes in \eqref{eq:coefficient-two-cross-21} therefore inflate
to zero on \(G\). Naturality of the five-term
sequence for \(1\to N\to G\to A\to1\) says that
the transgressions of the restrictions of \eqref{eq:coefficient-two-cross-22} to
\(N\) are exactly the pullbacks of \eqref{eq:coefficient-two-cross-21} to \(A\).
There is no additional affine component: this
statement comes from the displayed morphism of
extensions and its torus character.

On the arithmetic subgroup, the transgression
\[
 H^1(G_{K_2},\mathbf F_2)^H\longrightarrow H^2(A,\mathbf F_2)
\]
is injective, as in the complete arithmetic target
theorem. Thus \eqref{eq:coefficient-two-cross-21} identifies the restrictions as
\[
           \xi_{2,p}|_{G_{K_2}}=\kappa_{2,p},
       \qquad \widetilde\psi_p|_{G_{K_2}}=\psi_p|_{G_{K_2}}.
\]
This proves Theorem~\ref{thm:coefficient-two-cross-cups}.

\begin{corollary}\label{cor:all-arithmetic-quadratic-targets-lift}
The arithmetic restriction map is surjective:
\[
           X_2=H^1(N,\mathbf F_2)^G
              \longrightarrow Y_2=H^1(G_{K_2},\mathbf F_2)^H.
\]
Consequently the established structural sequence becomes
\begin{equation}
 0\longrightarrow\operatorname{Hom}_{\mathrm{cont}}
                         (B,\mathbf F_2)
   \longrightarrow X_2\longrightarrow Y_2\longrightarrow0.
                                                        \label{eq:coefficient-two-cross-23}
\end{equation}
This does not determine the left-hand term.
\end{corollary}

\begin{proof}
The complete target basis consists of \(\eta_3|_{G_{K_2}}\),
the \(\psi_p|_{G_{K_2}}\), the odd-pair characters
\(\kappa_{\ell,q}\), and the characters
\(\kappa_{2,p}\). The first two kinds already lift,
the odd-pair construction lifts the third kind, and
Theorem~\ref{thm:coefficient-two-cross-cups} lifts the
last kind. Every target character is a finite sum of
these basis elements. The kernel identification in
\eqref{eq:coefficient-two-cross-23} is the established relative character theorem
for \(G=NH\).
\end{proof}

\begin{remark}
The construction does not assert an equality on all
of \(N\) between distinct lifts of the same arithmetic
character. For \(p=3\), the earlier cubic frame
character \(\tau+\psi_3\) and \(\xi_{2,3}\) have the
same arithmetic restriction. Their difference, if
considered, is a character of the relative term in
\eqref{eq:coefficient-two-cross-23}; no vanishing or nonvanishing of that difference
is required here. Likewise, the companion character
\(\widetilde\psi_p\) is identified with \(\psi_p\)
arithmetically, without assuming their equality
on the full group.
\end{remark}
